\documentclass[11pt, a4paper]{article}

\usepackage[T1]{fontenc}
\usepackage[utf8]{inputenc}
\usepackage{tgpagella}
\usepackage[spanish, es-noshorthands]{babel}
\usepackage{amsmath, amssymb}
\usepackage[most]{tcolorbox}
\usepackage[margin=2.5cm]{geometry}
\usepackage{fancyhdr}

\pagestyle{fancy}
\fancyhf{}
\setlength{\headheight}{16pt}
\lhead{\textbf{M.Sc. Ing. Bernardo Quiroga Turdera}}
\rhead{IMT-221 | Instructor Trajectory Note}
\renewcommand{\headrulewidth}{0.4pt}
\definecolor{ucbblue}{RGB}{0, 51, 102}

\begin{document}

\begin{tcolorbox}[colback=black!5, colframe=black, title=\textbf{INTERNAL REGISTRY: EC3/EC4 Semester Trajectory Map}]
    Estado operacional al W10S03 de acuerdo a content-spec.
\end{tcolorbox}

\section*{1. Estado Curricular Inmediato}
\begin{itemize}
    \item \textbf{ALREADY CLOSED:} Ciclos escritos de U1 y U2. Fases de evaluación y recuperación teórica completadas.
    \item \textbf{PARALLEL CLOSURE:} Revisión de informes Hito 2, consolidación de calificaciones, reconciliación examen-reporte (cuando aplique) y microdefensas condicionales. \textbf{No deben consumir el bloque central en clase.}
    \item \textbf{COMPRESSED CARRY-FORWARD:} Conceptos de etapas A/B/AB (puente analógico). Desarrollado conceptualmente. Las derivaciones extensas matemáticas de eficiencia y sesgo (bias network design) se mantienen comprimidas / asíncronas para priorizar W10S03.
\end{itemize}

\section*{2. Main Forward Track (Unidad 3: Electrónica de Potencia)}
Trayectoria inamovible para las siguientes sesiones de clase (EC3):
\begin{enumerate}
    \item Modelo de conmutación MOSFET (Fundamentos entregados en W10S03).
    \item Selección de dispositivo desde un datasheet real de curso (Pendiente de selección paramétrica oficial).
    \item Dinámica de Gate Drive.
    \item Comportamiento frente a cargas resistivas e inductivas (Diodos Flyback).
    \item Evaluación térmica y modelo de pérdida de switching (Switching Loss).
    \item Topologías Half-Bridge / H-Bridge y prevención de Shoot-Through.
\end{enumerate}

\section*{3. EC4 - Protected Time (Unidad 4: Estabilidad y Feedback)}
\begin{tcolorbox}[colback=red!5, colframe=red!80!black, title=\textbf{RULE: TIMELINE PROTECTION}]
    No se permite que el cierre administrativo del Hito 2 o el retraso asimilado en MOSFET consuman el tiempo técnico reservado para EC4. EC4 debe mantenerse intacto.
\end{tcolorbox}
Tópicos protegidos: Realimentación, Respuesta en frecuencia, FFT/CSV para hardware/software, diagramas de Bode, márgenes de estabilidad y compensación integrada.

\section*{4. Next-Session Decision Gate (Action Items)}
Antes de avanzar a W10S04/W11, se requiere la ejecución de las siguientes definiciones paramétricas técnicas:
\begin{itemize}
    \item[$\square$] \textbf{MOSFET device selection:} Seleccionar número de parte oficial en base a restricciones de laboratorio (Ej. Voltaje de bus, Corriente RMS de carga, driver existente). \textit{IRLZ44N es meramente referencial.}
    \item[$\square$] \textbf{Gate-driver design constraints:} Definir tipo de control de disparo para el puente en laboratorio.
    \item[$\square$] \textbf{Load profiling:} Mediciones planeadas para la conmutación (Equipamiento de laboratorio).
\end{itemize}
\textit{Aviso: No se asignan fechas exactas futuras en esta hoja hasta que el calendario institucional restante no se haya reconciliado oficialmente tras el proceso de recuperación del Hito 2.}

\end{document}
